\subsection{Basic-to-Ordinary Narrowing: Exact Trace and MGU Transports}
\label{sec:p12-equational-narrowing-transport}

The source's Corollary4.14 says that the proof of basic narrowing also
works for ordinary narrowing and that the same transformation sequence
can be used. Its composite Solve/Narrow presentation then suppresses
the intervening syntactic steps. The transports below make those two
dependencies explicit, without identifying the two state spaces or
borrowing a new convergence hypothesis for a stepwise equivalence.
\cite[\S\S4.3.1--4.4.1, pp.~496--499]{BaaderSnyder2001}

\begin{hbdefinition}{def:p12-en-variants}{The exact two state interfaces and LP direction}
Use the source states and six rules in
\cref{def:p12-eu-gbn-states,sheet:p12-equational-unification-narrowing}.
Write \(A=\sigma_S\) for the simultaneous substitution of solved S
and \(D_S=\operatorname{Dom}(A)\). In B, P is the retained skeleton;
in N it receives each elimination substitution. C contains literal
constraints, not equations modulo R. The comparison scope is always
the initial problem variables, retained externally even when an
equation deletion removes one from the current state.
Throughout the state and trace theorems, B and N mean these six
primitive transformations only. Reduce, Propagation, redex abstraction,
Blocking and other refinements are not implicitly included in their
reachable-state invariant.

Fix the same chosen family \(\mathcal L\) of fresh LP identities
in both systems. For the directed composite of source \S4.4.1 it is
the forward rules \(l\to r\in R\). The source's B display in
\S4.3.1 inherits G's restrictions, while G selects identities from
\(E\cup E^{-1}\); if that inverse-identity reading is used,
take that same family in both systems instead. The transport is
valid for either fixed family, not an identification of their trees.
In every LP use, u is nonvariable at the selected skeleton occurrence;
if l is nonvariable, roots match and its newly introduced matching
constraint first undergoes Decomposition. Variants are fresh for the
whole earlier computation, as in source footnote11.
\cite[\S4.1.1, pp.~488--489; \S4.3.1, p.~496]{BaaderSnyder2001}
\end{hbdefinition}

\begin{hblemma}{thm:p12-en-state-invariant}{proved=annex:p12.en.state-invariant}{The actual reachable solved-state and global factor invariant}
In every reachable nonfailure B or N state, S is solved,
\(D_S\cap\operatorname{Var}(C)=\varnothing\), and A fixes its own
range. At eliminating \(x=^{?}t\) with \(\delta=\{x\mapsto t\}\),
x is outside \(D_S\), t contains neither x nor a member of \(D_S\),
and the new solved substitution is exactly \(A\delta\) on all
variables. In N, also \(D_S\cap\operatorname{Var}(P)=\varnothing\),
so \(P A=P\). Every literal S-solution \(\rho\) satisfies
\(A\rho=\rho\) globally. Thus at a terminal state \(\rho\)
itself is a factor for the returned answer, including initial
variables no longer occurring in the current P or C.
\end{hblemma}

\begin{hbdefinition}{def:p12-en-compile}{Compilation of a basic state and retained occurrences}
For a reachable basic state define its ordinary representation by
\[
         \operatorname{Inst}(P;C;S)=(P A;C;S).
\]
This is an equation-occurrence multiset, not a set obtained by
collapsing distinct equations. A nonvariable occurrence at address p
of a basic term remains at the same address of its A-instance, with
the same root and arity. Positions strictly inside a substituted
variable image may exist only in the latter; they receive no basic
address. Keep the first-Decomposition obligation on the same new
matching constraint in both representations.
\end{hbdefinition}

\begin{hbtheorem}{thm:p12-en-trace-transport}{proved=annex:p12.en.trace-transport}{Every basic trace has an actual ordinary trace with the identical answer}
Every finite B trace starting from \(P_0;\varnothing;\varnothing\)
has an N trace using the same chosen LP identities, such that at
each corresponding state \(P_N=P_B\sigma_S\) and C,S are literally
the same. Each of the six source rules and each matching-first
Decomposition is transported. If the basic trace terminates
\(\varnothing;\varnothing;S\), the ordinary trace terminates
with exactly that S and exactly that returned substitution.
No termination or confluence of R is required for this trace theorem.
The source's optional N Constrain variant that transfers bare e
instead of \(e A\) has the same effect on every reachable N state,
since \(P A=P\).
It is one-directional: an N step inside a substituted image need
not have a corresponding B step.
\end{hbtheorem}

\begin{hbmethod}{def:p12-en-source-composites}{Concrete U-MGU blocks and their accumulator interface}
For Solve take the literal equation \(Q=\{e A\}\). For directed
Narrow at a selected nonvariable u using fresh \(l\to r\), take
\(Q=\{l A=^{?}u A\}\). In the latter case first perform its
mandatory matching Decomposition, then run the actual removing-
binding U helper on the matching constraints. Let its idempotent
output be \(\eta\). It fixes variables outside \(\operatorname{Var}(Q)\)
and introduces no variable; in particular it fixes \(D_S\).
The new solved bindings are \(S\eta\cup[\eta]\), and their
substitution is \(A\eta\), not \(\eta A\).

In B the remaining/replaced skeleton is left uninstantiated; in N
it receives \(\eta\). The source's printed old-constraint context
\(C\eta\) may retain literal identities that a concrete set-based
helper run already deleted when a new constraint duplicated an old
one. Compare those presentations after Trivial deletion, never by
reclassifying C as equational. At the C-empty block boundaries
intended in \S4.4.1, this distinction is absent and the two composite
outputs are literal implementations of the registered rules.
\cite[\S2.2.3--2.2.4, pp.~454--457; \S4.4.1, p.~499]{BaaderSnyder2001}
\end{hbmethod}

\begin{hbtheorem}{thm:p12-en-mgu-blocks}{proved=annex:p12.en.mgu-blocks}{Finite composite blocks preserve the exact local factor condition}
These blocks terminate. A successful output has solved substitution
\(A\eta\). For every literal solution of old S and Q,
\(A\eta\rho=\rho\) globally. The output's literal S and old
constraint context describe exactly the old literal S,C,Q solutions,
with the stated harmless deletion of literal identities. A failing
match proves absence of that branch's literal witness, not absence
of all equational solutions of the original problem. Every successful
Solve or Narrow block is backwards equationally sound for its
selected problem equation.
\end{hbtheorem}

\begin{hblemma}{thm:p12-en-actual-lift}{proved=annex:p12.en.actual-lift}{A chosen innermost reduction lifts through an actual composite MGU}
Let \(\rho\) be a finite ground R-reduced witness for a reachable
B state, R-solving P and literally solving C,S. Suppose a selected term of P under
\(\rho\) has an innermost forward redex \(l\gamma\), where R
has ordinary rules (nonvariable left sides and no right-only variables).
Choose the least indexed rule matching that redex. Then its address
is nonvariable in the basic skeleton. The fresh matching substitution
\(\gamma\) extends \(\rho\) by normal proper-subterm images;
the literal matching constraint has that extension as a solution.
The actual U block therefore succeeds, with an explicit factor of
its \(\eta\), and preserves all old witness values. Under that
factor the new skeleton term is exactly the selected ground reduct.
The ordinary compiler transports the same block and factor. This
local lemma does not assume that every needed global witness is
groundable and normalizable in the original signature. Its groundness
and reduction tests are relative to the specified witness signature;
no ground confluence is needed for this one chosen-step construction.
\end{hblemma}

\begin{hbexample}{def:p12-en-transport-example}{Elimination changes the ordinary skeleton, not the basic occurrence map}
Let \(R=\{g(a)\to b\}\) and start with
\(\{h(x)=^{?}h(b),\ x=^{?}g(a)\}\). Constrain the second
equation and eliminate x. The B skeleton is
\(\{h(x)=^{?}h(b)\}\), the N skeleton is
\(\{h(g(a))=^{?}h(b)\}\), and both have S with
\(x\approx g(a)\). This is exactly \(P_N=P_B\sigma_S\).
N may narrow the g(a) child; B has a variable at its corresponding
address and may not use that occurrence. The trace theorem does not
transport this N-only step backwards or assert equality of search
trees. A different basic branch first rewrites the g(a) occurrence
while it is still in its original problem skeleton.
\end{hbexample}

\begin{hbexample}{def:p12-en-inverse-branch}{The weak-domain directed failure has a successful inverse-basic trace}
Use R and \(g(x)=^{?}a\) from
\cref{def:p12-eu-ground-domain-failure}, but now permit the inverse
identity of \(f(x,x)\to g(x)\). Choose fresh
\(g(v)\approx f(v,v)\) at the original g-root. The matching
constraint first decomposes to \(v=^{?}x\); Basic Elimination
sets v=x while the retained skeleton becomes \(f(v,v)=^{?}a\).
Its instantiated representation is \(f(x,x)=^{?}a\).
Next choose the forward identity \(f(p,q)\approx a\), fresh p,q.
Its matching constraint is \(f(p,q)=^{?}f(x,x)\); first
Decomposition and the literal solver give p=x,q=x. The remaining
skeleton is a=a, which Solve deletes. The terminal S is
\(\{v\approx x,p\approx x,q\approx x\}\), whose answer on
the initial scope \(\{x\}\) is identity. Both LP positions were
roots of actual basic skeletons. The trace compiler gives an N
trace in the same inverse-identity family with the identical answer.
It does not give a trace in the smaller directed-forward family.
Thus the directed weak-domain counterexample cannot be transferred
to the raw inverse-basic interpretation merely by calling both
procedures narrowing.
\end{hbexample}

\begin{hbtheorem}{thm:p12-en-terminal-normalization}{proved=annex:p12.en.terminal-normalization}{Terminal normalization preserves the answer class, not an active literal state}
For a terminal answer A and an ordinary R, if each image of A has
a finite reduction to a chosen normal form, normalize those images
to obtain \(\widehat A\), identity outside the same domain. Then
\(A=_R\widehat A\) pointwise, \(\widehat A\) is still idempotent,
and the two answers cover exactly the same targets modulo R on
the original scope. Soundness and completeness of any returned
family are preserved by this postprocessing. Termination of the
image reductions suffices for this assertion; confluence makes
the chosen forms unique. It does not license replacing a binding
in active S without transporting its literal constraints.
\end{hbtheorem}

\begin{hbexample}{def:p12-en-active-normalization-failure}{Equationally equal bindings can invalidate the literal active witness}
For \(R=\{f(a)\to a\}\), take the active state
\(P=\varnothing\), \(C=\{z=^{?}f(a)\}\),
\(S=\{x\approx f(a)\}\). Its literal witness gives
x=z=f(a). Normalizing only the x-binding to a removes that
witness, since the new S would require x=a literally. E-equality
between f(a) and a does not repair a literal S-equation. At a
terminal state with no C, normalized x=a is a sound representative
of the same equational answer class. These are different operations.
\end{hbexample}

\begin{hbtheorem}{thm:p12-en-infinite-search}{proved=annex:p12.en.infinite-search}{A finite convergent presentation can have an infinite basic and ordinary branch}
Let \(R=\{f(s(x))\to f(x),\ f(0)\to a\}\), with distinct
constants0,a and unary f,s. R terminates and is confluent on all
finite terms. Starting from \(f(y)=^{?}a\), there are successful
source composite B and N traces returning \(y\mapsto s^k(0)\)
for every \(k\geq0\), and an infinite branch that repeatedly
chooses the first rule. Restricted to the initial scope \(\{y\}\),
these answers form an infinite minimal complete set for that problem,
not just infinitely many redundant traces. Thus neither finite convergent R nor finite
matching U blocks imply termination of the nondeterministic narrowing
search. The repeated B steps occur in fresh replacement skeletons,
not inside previously bound images.
\end{hbtheorem}

\begin{proofboundary}{def:p12-en-completeness-boundary}{The transported component and the remaining universal dependency}
The trace theorem supplies the actual source Corollary4.14 transport:
every constructed successful basic witness trace is an ordinary trace
with the same answer and comparison factor. The MGU and lifting
lemmas supply the concrete \S4.4.1 composite blocks used on that
witness branch. They add no universal completeness premise and
do not replace the existing B lifting argument by a named result.

Source Theorem4.13 first grounds and reduces an arbitrary target.
Convergence only on the original signature's closed terms does not
establish that this operation preserves arbitrary parameter witnesses.
Example~\ref{def:p12-eu-ground-domain-failure} gives the precise
directed failure; Theorem~\ref{thm:p12-eu-bn-completeness} uses the
separately declared strengthened domain. The transport theorem does
not assert the missing grounding dependency under weaker original-
signature convergence. For the inverse-identity
reading of raw source B, a different full basic completeness argument
would have to establish its own witness normalization and protected-
position transport; the stepwise compiler alone does not do so.

These trace and factor results do not identify Fay's historical
presentation with normalization-based algorithms or establish the
arbitrary-theory hereditary refinements. In particular a run that
rewrites active S modulo R is not the ordinary N state transport
proved here.
\cite[Theorem4.13, Corollary4.14 and \S4.4.1, pp.~497--499]{BaaderSnyder2001}
\end{proofboundary}
